% =======================================================================================
\section{Conditional Flow Matching from Prior to Posterior}
\label{app:fm}
% =======================================================================================

\subsection{Objective}
\label{app:fm-objective}

For a window $w$ with data $d$ and conditioning input $y = y(d, w)$ (\cref{app:conditioning}),
the target is the posterior $p(X\mid y, w)\propto p(d\mid X)\,\pi_w(X)$ over the set of sources
$X\in\manifold^S$ in flow coordinates. We learn a time-dependent vector field
$\vnet(X, t; y, w)$ whose flow transports the prior to the posterior: draws
$X(0)\sim\pi_w$, integrated along
\begin{equation}
  \frac{\dd X}{\dd t} = \dot s(t)\,\vnet\bigl(X, t; y, w\bigr),\qquad t\in[0,1],
  \label{eq:ode}
\end{equation}
should end at $X(1)\sim p(\cdot\mid y,w)$. The clock $s(t)$ is the identity in plain flow
matching; \cref{app:fm-warp} warps it.

Training uses the conditional flow-matching construction
\citep{lipman2023flow,liu2023flow,albergo2023building} on a Riemannian manifold
\citep{chen2024flow}. A training example draws a window $w$, a target set
$X_1\sim\pi_w$ and the data $d\sim p(d\mid X_1)$ from the simulator, and an independent base
set $X_0\sim\pi_w$ in the same window. The conditional path is the geodesic between the
aligned sets of \cref{eq:set-log},
\begin{equation}
  X_s = \Exp_{X_0}\bigl(s\,\Log_{X_0}(X_1)\bigr),\qquad
  u_s = \frac{\dd X_s}{\dd s},
  \label{eq:cond-path}
\end{equation}
and the network regresses the conditional velocity,
\begin{equation}
  \mathcal L_{\rm flow}(\eta) =
  \E_{t\sim\Unif(0,1)}\;\E_{w,\,X_0,\,X_1,\,d}\;
  \frac{1}{11S}\Bigl\lVert \vnet\bigl(X_{s(t)}, t; y, w\bigr) - u_{s(t)}\Bigr\rVert^2 .
  \label{eq:loss-flow}
\end{equation}
Its minimiser is the marginal field $v^\star(X,s;y,w) = \E[u_s\mid X_s = X, y, w]$, which
generates the marginal path from $p_0 = \pi_w$ (because $X_0$ is independent of the data) to
$p_1 = p(\cdot\mid y,w)$ \citep{lipman2023flow,chen2024flow}. The base distribution is thus the
prior itself, not a Gaussian: on most factors of $\manifold$ the prior is uniform
(\cref{app:geometry-chart}), and starting from it means that a source the data say nothing
about only needs to stay where it is. Flow-matching posterior estimation with a Gaussian base
is studied by \citet{dax2023flow}.

\subsection{A Warped Clock}
\label{app:fm-warp}

The fine structure of a narrow posterior is resolved only at the very end of the path
(\cref{app:sensitivity}), where a uniform clock spends almost none of its samples. We keep
$t\sim\Unif(0,1)$ and the network's time input $t$, but place the point at
\begin{equation}
  s(t) = 1 - (1-t)^p,\qquad p = 3,
  \label{eq:warp}
\end{equation}
along the path, and keep the target $u_s = \dd X_s/\dd s$, which stays of order one as
$s\to1$. Equivalently, the regression in \cref{eq:loss-flow} is weighted by the density of
path positions
\begin{equation}
  q_p(s) = \frac{1}{p}(1-s)^{1/p-1},\qquad
  \Pr\bigl(1-s<\epsilon\bigr) = \epsilon^{1/p}.
  \label{eq:warp-density}
\end{equation}
For $p=3$, 21\% of the samples fall in the first half of the path ($1-2^{-1/3}$, against 50\%),
and 1\% of them within $10^{-6}$ of its end (against $10^{-6}$). The minimiser is unchanged,
since only the weighting over $s$ moves, and the sampler restores the clock through the factor
$\dot s(t) = p(1-t)^{p-1}$ in \cref{eq:ode}. Conditioning the network on $t$ rather than $s$
stretches the end of the path: $1-s = 10^{-6}$ is $1-t = 10^{-2}$, which the sinusoidal time
embedding resolves (\cref{app:architecture}). This plays the role of the non-uniform noise
levels used to train diffusion models \citep{karras2022elucidating,kingma2023understanding,esser2024scaling}, applied here to the end of the path rather than to its middle.

\subsection{Auxiliary Heads and the Full Loss}
\label{app:fm-aux}

The network has two further heads (\cref{app:architecture}): an estimate $\hat X_1$ of the
endpoint and a reconstruction $\hat y$ of the noiseless conditioning input, with losses
\begin{equation}
  \mathcal L_{x} = \E\,\frac{1}{11S}\bigl\lVert\Log_{X_1}(\hat X_1)\bigr\rVert^2,
  \qquad
  \mathcal L_{y} = \E\,\frac{1}{\lvert y\rvert}\bigl\lVert\hat y - y_{\rm clean}\bigr\rVert^2,
  \label{eq:aux-x}
\end{equation}
where $\Log$ is the set-aligned logarithm of \cref{eq:set-log} and $y_{\rm clean}$ is the
conditioning input computed from the noiseless signal. The training loss is
\begin{equation}
  \mathcal L = \frac{\mathcal L_{\rm flow}}{V_{\rm flow}}
  + a(\tau)\left[\frac{\mathcal L_{x}}{V_{x}} + \frac{\mathcal L_{y}}{V_{y}}\right],
  \qquad
  a(\tau) = \tfrac12 + \tfrac12\cos\bigl(\pi\min(\tau/\tau_a, 1)\bigr),
  \label{eq:loss}
\end{equation}
with $V_\bullet$ the running variance of each target over all training samples so far,
pooled over coordinates (a Welford accumulator), and $a(\tau)$ a cosine decay over the first
fraction $f_{\rm aux} = \tau_a/\tau_{\rm total}$ of training ($\tau$ counts epochs of 1000
steps). The original runs used $f_{\rm aux} = 0.5$; the recipe uses $f_{\rm aux}=0.1$
(\cref{app:xs}). Once $a=0$ the auxiliary heads get no gradient, their losses drift
(\cref{app:xs-dynamics}), and sampling never uses them. The logged losses are the raw mean
squared errors before the variance normalisation.

\subsection{Sampling}
\label{app:fm-sampling}

To sample a window we draw $N$ base sets from the prior, integrate \cref{eq:ode} with the
classical fourth-order Runge--Kutta scheme on 32 uniform steps in $t$, and map back with
$\varphi_w^{-1}$ (\cref{alg:sample}). The four stages are evaluated in the ambient
coordinates and only the full step is retracted onto the manifold with $\Exp$. Under the
warped clock, steps uniform in $t$ crowd towards the end of the path: the last of 32 ends at
$1-s = 32^{-3} = 3\times10^{-5}$. Each draw costs $4\times32 = 128$ network evaluations. Draws
are independent, so we push them through in chunks of 256 to bound the memory on B. The
integrator is converged at 32 steps: 64 and 128 steps change neither the widths nor the
detections (\cref{tab:numerics}).

\begin{algorithm}[h]
  \caption{One training step (per example; batched over 256 windows).}
  \label{alg:train}
  \begin{algorithmic}
    \STATE {\bfseries Input:} weights $\eta$, auxiliary weight $a$, running variances $V_\bullet$
    \STATE draw a window $w$ (\cref{app:signal-windows}) and $\theta_{1:S}\sim\pi_w$
    \STATE $\tilde d \leftarrow \sum_s\tilde h(\theta_s) + \tilde n$ \hfill (\cref{eq:data})
    \STATE $y\leftarrow y(\tilde d, w)$,\ \ $y_{\rm clean}\leftarrow y(\tilde h, w)$ \hfill (\cref{app:conditioning})
    \STATE $X_1\leftarrow\varphi_w(\theta_{1:S})$,\ \ $X_0\leftarrow\varphi_w(\theta'_{1:S})$, $\theta'_{1:S}\sim\pi_w$
    \STATE $t\sim\Unif(0,1)$,\ \ $s\leftarrow1-(1-t)^p$
    \STATE $X_s\leftarrow\Exp_{X_0}\bigl(s\Log_{X_0}(X_1)\bigr)$,\ \ $u_s\leftarrow\dd X_s/\dd s$ \hfill (\cref{eq:set-log,eq:cond-path})
    \STATE $(\hat v, \hat X_1, \hat y)\leftarrow\text{network}_\eta(X_s, t, y, \fw)$
    \STATE update $V_\bullet$; take a gradient step on \cref{eq:loss} (Muon/Adam, \cref{app:training})
  \end{algorithmic}
\end{algorithm}

\begin{algorithm}[h]
  \caption{Sampling the posterior of one window.}
  \label{alg:sample}
  \begin{algorithmic}
    \STATE {\bfseries Input:} data $\tilde d$ of window $w$, number of draws $N$, steps $K_t=32$
    \STATE $y\leftarrow y(\tilde d, w)$;\ \ $\Delta t\leftarrow1/K_t$
    \FOR{$n=1$ {\bfseries to} $N$ (in parallel)}
      \STATE $X\leftarrow\varphi_w(\theta_{1:S})$, $\theta_{1:S}\sim\pi_w$
      \FOR{$j=0$ {\bfseries to} $K_t-1$}
        \STATE $t\leftarrow j\Delta t$;\ \ $F(X,t) := p(1-t)^{p-1}\,\vnet(X,t;y,\fw)$
        \STATE $k_1\leftarrow F(X,t)$,\ \ $k_2\leftarrow F(X+\tfrac{\Delta t}{2}k_1, t+\tfrac{\Delta t}{2})$
        \STATE $k_3\leftarrow F(X+\tfrac{\Delta t}{2}k_2, t+\tfrac{\Delta t}{2})$,\ \ $k_4\leftarrow F(X+\Delta t\,k_3, t+\Delta t)$
        \STATE $X\leftarrow\Exp_X\bigl(\tfrac{\Delta t}{6}(k_1+2k_2+2k_3+k_4)\bigr)$
      \ENDFOR
      \STATE {\bfseries output} $\varphi_w^{-1}(X)$
    \ENDFOR
  \end{algorithmic}
\end{algorithm}
